% 09-statements.tex — Clause 9: Statements
\chapter{Statements}
\label{sec:09-statements}
\index{statement}

\pnum
A statement is a unit of execution.  Statements are terminated by line
feeds (\S\,\ref{sec:lex.whitespace}).  The grammar production for
\gramref{statement} is in Annex~A.

\section{Block statement}
\label{sec:stmt.block}
\index{block statement}

\pnum
A block statement \lain{\{ stmts... \}} creates a new scope
(\S\,\ref{sec:basics.scopes}) and executes each contained statement in
order.  Names declared inside a block are not accessible outside it.

\section{Immutable declaration}
\label{sec:stmt.decl}
\index{immutable declaration}

\pnum
\lain{name [Type] = expr} declares and initializes an immutable binding.
The type may be omitted (inferred from the initializer).  The binding
cannot be reassigned after initialization.

\begin{laincode}
x i32 = 42          // explicit type
y = 10              // type inferred (i32)
p Point = Point(1, 2)
\end{laincode}

\begin{constraint}
  An immutable declaration requires an initializer expression.
  The initializer shall not be \lain{undefined} (immutable bindings must
  have a definite value).
\end{constraint}

\begin{constraint}
  When a type annotation is present, the type of the initializer shall be
  compatible with the declared type (\S\,\ref{sec:types.compat}).
\end{constraint}

\begin{note}
  When the type is omitted, the distinction between declaration and
  reassignment is resolved by scope: \lain{x = 5} where \lain{x} is not
  yet declared creates a new immutable binding; where \lain{x} is already
  declared as \lain{var}, it is a reassignment.
\end{note}

\section{Mutable variable declaration}
\label{sec:stmt.var}
\index{variable declaration}

\pnum
\lain{var name Type = expr} declares a mutable variable.  Mutable variables
can be reassigned.  The type may be omitted if it can be inferred from the
initializer.

\begin{laincode}
var counter i32 = 0
counter = counter + 1       // OK: var is mutable

var buf u8[256]              // OK: type present, defaults to undefined
\end{laincode}

\begin{constraint}
  A mutable declaration shall include either an explicit type annotation
  or an initializer expression.  A declaration with neither is
  \illformed{}~(\diagcode{E100}).
\end{constraint}

\begin{constraint}
  When both a type annotation and an initializer are present, the type of
  the initializer shall be compatible with the declared type
  (\S\,\ref{sec:types.compat}).
\end{constraint}

\section{Assignment statement}
\label{sec:stmt.assign}
\index{assignment statement}

\pnum
An assignment statement \lain{lvalue = expr} stores the value of
\textit{expr} into the storage designated by \textit{lvalue}.
An assignment target without a type annotation is always a reassignment,
never a declaration.

\begin{constraint}
  The left-hand side shall be a mutable lvalue (\S\,\ref{sec:basics.lvalues}).
  Assigning to an immutable binding is \illformed{}~(\diagcode{E009}).
\end{constraint}

\begin{constraint}
  Overwriting a linear variable whose current value has not been consumed
  would leak that value.  Such an assignment is
  \illformed{}~(\diagcode{E003}).
\end{constraint}

\begin{note}
  Assigning to a linear variable after its previous value has been consumed
  (e.g., loop reassign-and-consume pattern) is permitted.  The implementation
  resets the variable's linear state to Unconsumed after a legal overwrite.
\end{note}

\begin{example}
\begin{laincode}
import std.fs

proc main() int {
    var f = open_file("data.txt", "r")
    close_file(mov f)               // f is now Consumed
    f = open_file("data.txt", "r")  // OK: overwriting Consumed state
    close_file(mov f)
    return 0
}
\end{laincode}
\end{example}

\section{Expression statement}
\label{sec:stmt.expr}
\index{expression statement}

\pnum
An expression followed by a statement terminator is an expression
statement.  The value of the expression is discarded.

\section{If statement}
\label{sec:stmt.if}
\index{if statement}

\pnum
\lain{if cond \{ body \} [else if cond \{ body \}]... [else \{ body \}]}
evaluates \textit{cond}; if true, executes the then-body; otherwise,
tests the \lain{else if} clauses in order; if none matches, executes the
\lain{else} body (if present).

\begin{constraint}
  \textit{cond} shall have type \lain{bool}.
\end{constraint}

\pnum
\textbf{Range refinement}: inside the then-branch of \lain{if x < N \{ \}},
the value range analysis (\S\,\ref{sec:13-vra}) knows that
$x \in (-\infty, N-1]$; inside the else-branch, $x \in [N, +\infty)$.

\pnum
\textbf{In-guard}: when \textit{cond} contains or is an \lain{in}
expression (\S\,\ref{sec:expr.in}), the in-guard is active inside the
then-branch only.

\pnum
\textbf{Branch consistency}: if one branch consumes a linear variable,
all branches shall consume it.  If branches disagree on the linear state
of a variable from the enclosing scope, the program is
\illformed{}~(\diagcode{E003}).

\section{While statement}
\label{sec:stmt.while}
\index{while statement}

\pnum
\lain{while cond \{ body \}} repeatedly evaluates \textit{cond} and, if
true, executes \textit{body}.

\begin{constraint}
  An unbounded \lain{while} loop (one without a \lain{decreasing} measure)
  shall not appear inside a \lain{func}.  Such use is
  \illformed{}~(\diagcode{E011}).
\end{constraint}

\begin{constraint}
  \textit{cond} shall have type \lain{bool}.
\end{constraint}

\pnum
\textbf{In-guard}: when \textit{cond} contains or is an \lain{in}
expression, the in-guard is active inside the loop body for each
iteration.

\subsection{Bounded while statement}
\label{sec:stmt.while.bounded}
\index{bounded while loop}
\index{decreasing measure}

\pnum
\lain{while cond decreasing M \{ body \}} is a bounded loop.  In addition
to the while semantics above, the implementation shall statically verify:

\begin{enumerate}
  \item \textbf{Non-negativity} (\diagcode{E080}): the implementation shall
        prove that $M \geq 0$ whenever \textit{cond} is true.
  \item \textbf{Variable dependency} (\diagcode{E081}): $M$ shall
        reference at least one variable that is assigned in \textit{body}.
        An $M$ that references no such variable would never decrease.
  \item \textbf{Strict decrease} (\diagcode{E082}): the implementation
        shall prove that every execution of \textit{body} strictly
        decreases the value of $M$.
\end{enumerate}

\pnum
A bounded \lain{while} loop is permitted inside a \lain{func} because the
implementation has statically verified termination.

\pnum
The following patterns are accepted by the standard verification algorithm:

\begin{longtable}{@{}lll@{}}
\toprule
\textbf{Condition} & \textbf{Measure} & \textbf{Why} \\
\midrule
\endfirsthead
\toprule
\textbf{Condition} & \textbf{Measure} & \textbf{Why} \\
\midrule
\endhead
\bottomrule
\endlastfoot
\lain{i < n}    & \lain{n - i}      & $i$ increases $\Rightarrow$ measure decreases \\
\lain{n > 0}    & \lain{n}          & $n$ decreases directly \\
\lain{a < b}    & \lain{b - a}      & either $a$ increases or $b$ decreases \\
\lain{i in arr} & \lain{arr.len - i}& \lain{in} implies $i < \texttt{arr.len}$ \\
\end{longtable}

\begin{example}
\begin{laincode}
func count(n int) int {
    var i = 0
    while i < n decreasing n - i {
        i = i + 1
    }
    return i
}
\end{laincode}
\end{example}

\begin{rationale}
  The \lain{decreasing} keyword bridges the gap between pure functions
  (which must terminate) and the common pattern of finite-input state
  machines (lexers, parsers) that are provably terminating.  The programmer
  states the termination argument; the compiler verifies it.
\end{rationale}

\section{Return statement}
\label{sec:stmt.return}
\index{return statement}

\pnum
\lain{return [expr]} terminates execution of the current function or
procedure and transfers the value of \textit{expr} (if present) to the
caller.

\begin{constraint}
  At the point of a \lain{return}, all linear variables in scope shall
  have been consumed.  A live (Unconsumed) linear variable at \lain{return}
  is \illformed{}~(\diagcode{E003}).
\end{constraint}

\begin{constraint}
  \textbf{Q-009.}
  A return statement shall not be prefixed with \lain{mov} or \lain{var}.
  The forms \lain{return mov x} and \lain{return var x} are
  syntactically \illformed{}~(\diagcode{E100}).  Ownership transfer in
  return position is implicit from the function's declared return type.
\end{constraint}

\section{Defer statement}
\label{sec:stmt.defer}
\index{defer statement}

\pnum
\lain{defer \{ body \}} schedules \textit{body} for execution immediately
before the enclosing scope exits, on all exit paths (\lain{return},
\lain{break}, \lain{continue}, or fall-through).  Multiple \lain{defer}
statements execute in LIFO order.

\pnum
\textbf{Interaction with linear types}: if a linear variable is consumed
inside a \lain{defer} body, the variable is marked as defer-consumed.  It
is treated as unconsumed for the purposes of the main body (it may be used
before the scope exit) and as consumed at scope exit.

\begin{constraint}
  A \lain{defer} body shall not contain \lain{return}, \lain{break}, or
  \lain{continue} statements that would transfer control out of the deferred
  body.
\end{constraint}

\begin{example}
\begin{laincode}
import std.fs

proc process() int {
    var f = open_file("data.txt", "r")
    defer { close_file(mov f) }   // f consumed at scope exit
    // ... use f ...
    return 0
}
\end{laincode}
\end{example}

\section{Unsafe block}
\label{sec:stmt.unsafe}
\index{unsafe block}

\pnum
\lain{unsafe \{ body \}} is an unsafe block (see \S\,\ref{sec:18-unsafe}).
It creates a scope in which certain static checks are relaxed.

\section{Comptime if statement}
\label{sec:stmt.comptime-if}
\index{comptime if}

\pnum
\lain{comptime if cond \{ then \} [else \{ else \}]} evaluates
\textit{cond} at compile time.  Only the taken branch is type-checked and
emitted; the other branch is discarded.  \textit{cond} shall be a
compile-time constant expression (\S\,\ref{sec:14-generics}).

\begin{note}
  \lain{comptime if} is the Lain equivalent of C preprocessor
  \texttt{\#if}: it enables platform-specific or type-specific code
  selection without dead-code warnings.
\end{note}

\section{For statement}
\label{sec:stmt.for}
\index{for statement}

\pnum
\lain{for i in start..end \{ body \}} iterates with \textit{i} ranging
from \textit{start} (inclusive) to \textit{end} (exclusive).  The loop
variable \textit{i} is immutable; its value range is statically known to
be $[\textit{start}, \textit{end}-1]$ throughout the body, enabling safe
array indexing without runtime bounds checks.

\pnum
\lain{for} loops are permitted in both \lain{func} and \lain{proc}.

\begin{constraint}
  \textit{start} and \textit{end} shall have integer types.
\end{constraint}
